(a) For the node, the chart $y=xu$ of the strict transform has equation $u-x^4(1+u^6)=0$. Above the origin there is one point $(x,u)=(0,0)$, where the derivative with respect to $u$ is $1$. The other chart gives a second nonsingular point by symmetry. The blowup is an isomorphism away from the origin, which is the only singular point of the original curve.

For the cusp, the same chart has equation $x-u^2-x^2(1+u^4)=0$. Its exceptional fiber is the single point $(0,0)$, where the derivative with respect to $x$ is $1$. The chart $x=ty$ has equation $t^3y-1-y^2(t^4+1)=0$, so it has no exceptional point. Consequently the blowup is nonsingular in characteristic $0$, and more generally when the characteristic is different from $2,7,13$. In characteristics $7$ and $13$, the two additional singularities found in \exref{I.5.1} remain, since they are away from the center of the blowup.

(b) Translate $P$ to the origin and make a linear change of coordinates so that the equation is $xy+f_3+f_4+\cdots=0$. In the chart $y=xu$, the strict transform has equation $u+x f_3(1,u)+x^2f_4(1,u)+\cdots=0$. Its only exceptional point is $(0,0)$, where the derivative with respect to $u$ is $1$. The chart $x=ty$ gives the other nonsingular point. They correspond to the distinct tangent directions $y=0$ and $x=0$.

(c) In the chart $x=ty$, the strict transform is $t^2-y^2(t^4+1)=0$. The exceptional fiber is $(t,y)=(0,0)$, and its tangent cone is $(t-y)(t+y)=0$, so this point is a node. The other chart has equation $1-x^2(1+u^4)=0$ and no exceptional point. Part (b) resolves the remaining node.

(d) The lowest term of $y^3-x^5$ has degree $3$. In the chart $y=xu$, the strict transform is $u^3=x^2$, with a cusp of multiplicity $2$ at the origin. The chart $x=ty$ has equation $1=t^5y^2$ and no exceptional point. Blow up the cusp using $x=uv$; its strict transform is $u=v^2$, which is nonsingular. The other chart has no additional exceptional point, so the second blowup resolves the singularity.
